% Math graph Paper created for Caleb using the help of Gemini
%
% Eric Westervelt
% 12/27/2025

\documentclass[letterpaper]{article}
\usepackage[margin=0in]{geometry}
\usepackage{tikz}
\usetikzlibrary{calc}
\usepackage[pdftex, hidelinks]{hyperref}

\begin{document}
\pagestyle{empty}

% Define the default appearance for text fields
\def\DefaultHugeBold{
  \DeclareFontFamily{T1}{cmr}{}
  \DeclareFontShape{T1}{cmr}{bx}{n}{<-> cmbx12}{}
  \fontsize{24}{28}\selectfont\bfseries
}

% Command to draw a box number field
\newcommand{\BoxNum}[3]{% #1=X, #2=Y, #3=Value
    \node[anchor=north west] at (#1, #2) {
        \TextField[name=field#3, width=35pt, height=26pt, bordercolor=, align=1, value=#3, charsize=16pt]{}
    };
}

% Command to generate one page
% #1 = Starting Number, #2 = Header Mode (h for header, n for no header)
\newcommand{\ProblemPage}[2]{
    \begin{Form}
    \begin{tikzpicture}[remember picture, overlay, shift={(current page.south west)}]
        % 1. Create the fine graph paper
        \draw[cyan!30, thin] (0,0) grid [step=0.25in] (8.5in, 11in);

        % 2. Header and Grid Logic
        \if#2h
            % Page WITH Header (Pages 1 and 3)
            \node[anchor=west, font=\Large\bfseries] at (0.4in, 10.35in) {Lesson:};
            \node[anchor=west] at (1.2in, 10.45in) {\TextField[name=Lesson#1, width=0.7in, height=0.3in, bordercolor=]{}};
            \draw[ultra thick] (1.2in, 10.25in) -- (2in, 10.25in);

            \node[anchor=west, font=\Large\bfseries] at (2.1in, 10.35in) {Date:};
            \node[anchor=west] at (2.8in, 10.45in) {\TextField[name=Date#1, width=2in, height=0.3in, bordercolor=]{}};
            \draw[ultra thick] (2.8in, 10.25in) -- (4.9in, 10.25in);
            
            \node[anchor=west, font=\Large\bfseries] at (5in, 10.35in) {Name:};
            \node[anchor=west] at (5.75in, 10.45in) {\TextField[name=Name#1, width=2.35in, height=0.3in, bordercolor=]{}};
            \draw[ultra thick] (5.75in, 10.25in) -- (8.2in, 10.25in);

            \draw[black, ultra thick] (4.25in, 0in) -- (4.25in, 10in);
            \draw[black, ultra thick] (0, 10in) -- (8.5in, 10in);
            \draw[black, ultra thick] (0, 6.75in) -- (8.5in, 6.75in);
            \draw[black, ultra thick] (0, 3.5in) -- (8.5in, 3.5in);
            %\draw[black, ultra thick] (0, 0.05in) -- (8.5in, 0.05in);

            \def\yA{9.95in} \def\yB{6.7in} \def\yC{3.45in}
        \else
            % Page WITHOUT Header (Page 2)
            \draw[black, ultra thick] (4.25in, 0in) -- (4.25in, 11in);
            %\draw[black, ultra thick] (0, 11in) -- (8.5in, 11in);
            \draw[black, ultra thick] (0, 7.25in) -- (8.5in, 7.25in);
            \draw[black, ultra thick] (0, 3.75in) -- (8.5in, 3.75in);
	 %\draw[black, ultra thick] (0, 0in) -- (8.5in, 0in);

            \def\yA{10.8in} \def\yB{7.2in} \def\yC{3.7in}
        \fi
        
        % 3. Calculate Box Numbers
        \pgfmathtruncatemacro{\numA}{#1}
        \pgfmathtruncatemacro{\numB}{#1 + 2}
        \pgfmathtruncatemacro{\numC}{#1 + 4}
        \pgfmathtruncatemacro{\numD}{#1 + 6}
        \pgfmathtruncatemacro{\numE}{#1 + 8}
        \pgfmathtruncatemacro{\numF}{#1 + 10}

        % Row 1
        \BoxNum{0.3in}{\yA}{\numA}
        \BoxNum{4.4in}{\yA}{\numB}
        
        % Row 2
        % Conditional check for row 2 (For the 25-29 range)
        \ifnum\numC<30 \BoxNum{0.3in}{\yB}{\numC}\else\BoxNum{0.3in}{\yB}{}  \fi
        \ifnum\numD<30 \BoxNum{4.4in}{\yB}{\numD}\else\BoxNum{4.4in}{\yB}{} \fi
        
        % Row 3
        \ifnum\numE<30 \BoxNum{0.3in}{\yC}{\numE}\else\BoxNum{0.3in}{\yC}{} \fi
        \ifnum\numF<30 \BoxNum{4.4in}{\yC}{\numF}\else\BoxNum{4.4in}{\yC}{} \fi
        
    \end{tikzpicture}
    \end{Form}
    \clearpage
}

% --- Generate the pages ---
\ProblemPage{1}{h}  % Page 1: 1, 3, 5, 7, 9, 11
\ProblemPage{13}{n} % Page 2: 13, 15, 17, 19, 21, 23
\ProblemPage{25}{h} % Page 3: 25, 27, 29 (31 is excluded)

\end{document}